%% ma: L11-B03-0012
%% lop: 11
%% chuong: 1
%% bai: 3
%% dang: 11.3.8
%% loai: DS
%% muc: [TH, NB, TH, TH]
%% dap_an: "SSĐĐ"
%% nguon: "de-cuong-giua-ky-1-toan-11-nam-2026-2027-truong-thpt-hoang-van-thu-ha-noi-2-8.pdf (M08, Đề cương giữa kì 1 Toán 11 THPT Hoàng Văn Thụ, Hà Nội), Phần 2 Câu 3"
%% kien_thuc: "Bài 3 (hàm số lượng giác)"
%% trang_thai: chinh-thuc
%% ghi_chu: "Sửa lỗi chính tả \"tâm tương\" thành \"tâm trương\" ở ý d. Đã duyệt 29/10/2026."
\begin{ex}
Huyết áp là áp lực máu cần thiết tác động lên thành động mạch nhằm đưa máu đi nuôi dưỡng các mô trong cơ thể. Nhờ lực co bóp của tim và sức cản của động mạch mà huyết áp được tạo ra. Giả sử huyết áp của một người thay đổi theo thời gian được cho bởi công thức $p(t)=120+15\cos150\pi t$, trong đó $p(t)$ là huyết áp tính theo đơn vị mmHg (milimét thủy ngân) và thời gian $t$ tính theo đơn vị phút. Xét tính đúng sai của các khẳng định sau:
\choiceTF{Hàm số $p(t)$ tuần hoàn với chu kì $\dfrac{\pi}{75}$}{Tại thời điểm $t=0$, huyết áp của người này là $120$ mmHg}{\True Huyết áp tâm thu (huyết áp cao nhất) của người này là $135$ mmHg}{\True Huyết áp tâm trương (huyết áp thấp nhất) của người này là $105$ mmHg}
\loigiai{a) Sai: chu kì $T=\dfrac{2\pi}{150\pi}=\dfrac{1}{75}$ (phút).

b) Sai: $p(0)=120+15=135$ (mmHg).

c) Đúng: $p(t)\le120+15=135$, đạt khi $\cos150\pi t=1$.

d) Đúng: $p(t)\ge120-15=105$, đạt khi $\cos150\pi t=-1$.}
\end{ex}
